% Generated draft: requires mathematical and source review.
\documentclass[11pt]{article}
\usepackage[margin=25mm]{geometry}
\usepackage{fontspec}
\setmainfont{Latin Modern Roman}
\usepackage{amsmath,amssymb,mathtools,unicode-math}
\setmathfont{Latin Modern Math}
\usepackage{enumitem,needspace,xurl,hyperref}
\setlength{\parindent}{0pt}
\setlength{\parskip}{5pt}
\setlength{\emergencystretch}{3em}
\newcommand{\N}{\mathbb{N}}
\newcommand{\Z}{\mathbb{Z}}
\newcommand{\Q}{\mathbb{Q}}
\newcommand{\R}{\mathbb{R}}
\newcommand{\C}{\mathbb{C}}
\newcommand{\F}{\mathbb{F}}
\newcommand{\A}{\mathbb{A}}
\newcommand{\PP}{\mathbb{P}}
\newcommand{\E}{\mathbb{E}}
\newcommand{\eps}{\varepsilon}
\newcommand{\dd}{\,\mathrm{d}}
\begin{document}
\section*{UnsolvedMath ID: 30003416}
\textbf{Source problem code:} catalogue entry
\subsection{Definitions and mathematical statement}
The source does not define the family \((G_\xi)_{\xi<\varepsilon_0}\). Consequently, the following records its quantified assertion, but a complete mathematical formulation requires the intended definition of that family.

Under the conventional interpretation \(I=[0,1]\), let \(\mathrm{PL}_+(I)\) be the group, under composition, of increasing piecewise-linear homeomorphisms of \(I\) with finitely many linear pieces. Such homeomorphisms fix both endpoints. The source does not explicitly specify the interval or the convention concerning the number of pieces; these conventions require confirmation.

Let Thompson's group \(F\) be the subgroup of \(\mathrm{PL}_+([0,1])\) consisting of maps whose breakpoints are dyadic rationals and whose slopes belong to \(\{2^n:n\in\Z\}\). A group embedding is an injective group homomorphism, and containing a copy of \(F\) means having a subgroup isomorphic to \(F\). Let \(\varepsilon_0\) be the least nonzero ordinal \(\alpha\) satisfying \(\omega^\alpha=\alpha\). A group is elementary amenable if it belongs to the smallest class containing all finite groups and all abelian groups and closed under subgroups, quotients, group extensions, and directed unions.

The supplied conjecture asks whether
\[
\begin{gathered}
\text{for every finitely generated }G\leq\mathrm{PL}_+(I),\\
\bigl(\text{there is no injective homomorphism }F\to G\bigr)
\ \Longrightarrow\ 
\bigl(\exists\xi<\varepsilon_0\;\exists\text{ injective homomorphism }G\to G_\xi\bigr).
\end{gathered}
\]
The source also asks, as an asserted particular case, whether
\[
\begin{gathered}
\text{for every finitely generated elementary amenable }H\leq F,\\
\exists\xi<\varepsilon_0\;\exists\text{ injective homomorphism }H\to G_\xi.
\end{gathered}
\]
Here \(G_\xi\) must refer to the specific hierarchy intended by the source, rather than an arbitrary ordinal-indexed family. Its construction and indexing conventions are absent from the supplied record. The particular-case wording is preserved from the source; no proof of that implication is supplied here.
\subsection{Short English statement}
The conjecture asks whether every finitely generated group of piecewise-linear interval homeomorphisms containing no subgroup isomorphic to Thompson's group \(F\) embeds into a member of a specified hierarchy below \(\varepsilon_0\). It also asks this for finitely generated elementary amenable subgroups of \(F\). The hierarchy is undefined in the supplied record, so the statement requires source clarification and human mathematical review.
\subsection{Sources}
\begin{itemize}
\item Local UnsolvedMath registry, \texttt{data/unsolvedmath/problems.json}, permanent ID \(30003416\), source code \texttt{catalogue entry}, title: Embedding Interval Groups Without Thompson F into the G-Hierarchy. The supplied record belongs to the Oberwolfach Reports Open Problems collection and lists proposed year 2017; no proposer is named.
\item Supplied report citation: Set Theory (2018), Oberwolfach Reports, DOI 10.4171/owr/2017/11, pp.\ 527--589. \url{https://doi.org/10.4171/owr/2017/11}. The citation's year label is retained as supplied; its relation to the 2017 DOI component and proposed year requires confirmation.
\item Supplied literature reference: Groups of fast homeomorphisms. \url{https://research-repository.st-andrews.ac.uk/bitstream/10023/18600/1/fastgen_JCA_final.pdf}.
\item Supplied literature reference: Complexity among finitely generated subgroups of F. \url{https://arxiv.org/abs/1711.10998}.
\item Supplied literature reference: OWR Set Theory report. \url{https://ems.press/content/serial-article-files/46674}.
\item Literature assessment dated exactly \texttt{2026-08-21}: the supplied curation labels the problem open and states that the cited subgroup-classification paper presents the bi-embeddability and embedding assertions as conjectures, while its explicit families and structural theorems do not resolve the general claims. This is supplied dated curation derived from the OWR record, not an independently verified literature assessment. The references and their claims have not been independently checked here.
\item The registry labels the statement \texttt{corrected\_verified} and says that notation, typography, and sentence structure were repaired while preserving mathematical content. These are supplied editorial assertions, not mathematical verification of this draft.
\end{itemize}
\end{document}
